\slajdid{03-s019}
\begin{frame}[label=03-s019]
\gusttekst
Prilikom pravljenja mreža većih kapaciteta, možemo se poslužiti idejama koje smo imali kod prethodnih kola. Na primer da je potrebno napraviti koder prioriteta sa 16 ulaza
\begin{columns}[c,onlytextwidth]
\begin{column}{.43\textwidth}\centering
\begingroup
\def\DEenable{0}
% \KPSimbol{ime}{x}{y}; DEenable bira EI/EO varijantu; font 8 pt u pozivaocu.
\newcommand{\KPSimbol}[3]{%
\begin{scope}[shift={(#2,#3)}]
\draw[DEvod] (0,0) rectangle (2.1,1.65);
\ifnum\DEenable=1
 \def\DEgsheight{1.05}
 \node at (1.0,1.05) {KP};
 \node[anchor=north,inner sep=1pt] at (1.0,1.65) {EI};
 \draw[DEvod] (1.0,1.65)--(1.0,1.8);
 \coordinate (#1-ei) at (1.0,1.8);
 \node[anchor=south,inner sep=1pt] at (1.0,0) {EO};
 \draw[DEvod] (1.0,0)--(1.0,-.15);
 \coordinate (#1-eo) at (1.0,-.15);
 \foreach \n/\y in {1/.62,0/.25}{
  \node[anchor=east,inner sep=1pt] at (2.1,\y) {A\n};
  \draw[DEvod] (2.1,\y)--(2.3,\y);
  \coordinate (#1-a\n) at (2.3,\y);
 }
\else
 \def\DEgsheight{.95}
 \node at (1.0,1.43) {KP};
 \foreach \n/\x in {1/.7,0/1.45}{
  \node[anchor=south,inner sep=1pt] at (\x,0) {A\n};
  \draw[DEvod] (\x,0)--(\x,-.15);
  \coordinate (#1-a\n) at (\x,-.15);
 }
\fi
\foreach \n/\y in {3/1.2,2/.89,1/.58,0/.27}{
 \node[anchor=west,inner sep=1pt] at (0,\y) {I\n};
 \draw[DEvod] (-.2,\y)--(0,\y);
 \coordinate (#1-i\n) at (-.2,\y);
}
\node[anchor=east,inner sep=1pt] at (2.1,\DEgsheight) {GS};
\draw[DEvod] (2.1,\DEgsheight)--(2.3,\DEgsheight);
\coordinate (#1-gs) at (2.3,\DEgsheight);
\end{scope}}

\begin{tikzpicture}[x=.8cm,y=.8cm,font=\fontsize{8}{9}\selectfont]
\foreach \k/\y in {3/6,2/4,1/2,0/0}{
 \KPSimbol{p\k}{0}{\y}
 \foreach \n in {3,2,1,0}{
  \pgfmathtruncatemacro{\ii}{4*\k+\n}
  \draw[DEvod] (p\k-i\n)--++(-.15,0) node[left,inner sep=1pt] {I\ii};
 }
}
\KPSimbol{f}{3.5}{3}
\foreach \k/\x in {3/3.0,2/2.7,1/2.7,0/3.0}
 \draw[DEvod] (p\k-gs)--(\x,0 |- p\k-gs)|-(f-i\k);
\draw[DEvod] (f-gs)--++(.2,0) node[right] {GS};
\draw[DEvod] (f-a1)--++(0,-.2) node[below] {A3};
\draw[DEvod] (f-a0)--++(0,-.2) node[below] {A2};
\end{tikzpicture}
\endgroup

\end{column}
\begin{column}{.54\textwidth}
Sa slike je vidljiva ideja da se iskoristi osobina kodera prioriteta i u drugom nivou. Znači ako ima na bilo kojem koderu prioriteta u prvom nivou aktivnih logičkih jedinica, biće aktivan i njegov izlazni signal GS, pa će koder prioriteta u drugom nivou dati odgovarajući indeks kodera prioriteta prvog nivoa (određen po prioritetu GS signala), što su u stvari signali A3 i A2 za „veliki“ koder. Ostaje problem kako odrediti signale A1 i A0. \alert{Ovde su u pitanju izlazi pa njihovo direktno spajanje sigurno ne dolazi u obzir.} A i logički njihovo direktno spajanje ne bi donelo validnu informaciju. U tom smislu bi morali na osnovne kodere prioriteta dodati još neke ulaze i izlaze, kako bi omogućili laku identifikaciju koji signali A1 i A0 su validni, a i da omogućimo njihovo lako logičko spajanje.
\end{column}
\end{columns}
\end{frame}
